\section{本讲小结}
交流计算先区分瞬时值、平均值、有效值和相位。有效值按完整周期内的平方平均定义，非正弦波不能一律用峰值除以 $\sqrt2$。线性同频稳态使用阻抗相量；串联加电压、并联加电流，都要保留相位。二极管与整流过程先由导通边界分段，再积分求平均值、有效值和功率。

谐振对应电抗或电纳的抵消，电源电压、电源电流和支路量应分别辨认。功率因数补偿改变电源无功电流；理想电容不改变原负载在固定端压下的有功功率。理想变压器按端口参考方向连接电压、电流，并按匝数比平方折算阻抗。最大功率与最高效率分别由输出和损耗模型决定。

三相计算先分清相电压、线电压、相电流和线电流。中性点接通与浮置给出不同的负载端压；浮置时应先求中性点电势。输电计算使用实际线路电流、两端电压和有功功率，不能把含电抗的交流负载直接当作固定纯电阻。

位移电流通过变化的电通量连接充电区与导线区。集中 LC 电路描述局部储能交换，分布传输线同时包含空间推进与电、磁两种场能；无损线的行波阻抗连接电压和电流，不代表发热电阻。遇到反射或非均匀场，须重新辨认积分面、边界和传播方向。
